\section{A rational parametrization and its integral defect}\label{sec:scaling}
For the signature $(3,5,7)$, rational points are easy to manufacture. Primitive integer points are not. The following polynomial identity makes the difference exact.
\begin{proposition}[Universal power identity]\label{prop:universal}
In any commutative ring, put $s=1-t$ and
\[
A_0=t^{12}s^{28},\qquad B_0=t^7s^{17},\qquad C_0=t^5s^{12}.
\]
Then $A_0^3+B_0^5=C_0^7$. Weighted scaling
\begin{equation}\label{eq:weightedscale}
(A,B,C)=(\lambda^{35}A_0,\lambda^{21}B_0,\lambda^{15}C_0)
\end{equation}
preserves this equality.
\end{proposition}
\begin{proof}
The two terms on the left are $t^{36}s^{84}$ and $t^{35}s^{85}$. Their sum is $t^{35}s^{84}(t+s)=t^{35}s^{84}$, which is $C_0^7$. The three scaling weights multiply the powered terms by the same factor $\lambda^{105}$.
\end{proof}
Over a field, if $t$ and $1-t$ are nonzero, all three bases are nonzero and $A_0^3/C_0^7=t$. Thus every regular parameter is locally populated. This works modulo arbitrary prime powers whenever $t$ and $1-t$ are units. A finite field test that asks only whether such a nonzero local solution exists cannot remove any of these parameter values.

For a nonzero rational solution, define $t=A^3/C^7$ and $\lambda=C^8/(AB^4)$. Substitution recovers \eqref{eq:weightedscale}. If the original bases are positive, then $0<t<1$ and $\lambda>0$. The rational parametrization therefore captures the entire positive rational problem. Its missing condition is exactly primitive integrality.

\begin{theorem}[Primitive scaling criterion]\label{thm:scaling}
Fix a rational $0<t<1$. There is a positive rational $\lambda$ for which \eqref{eq:weightedscale} gives primitive positive integers if and only if, for every rational prime $\ell$,
\begin{align*}
v_\ell(t)>0&\quad\Longrightarrow\quad3\mid v_\ell(t),\\
v_\ell(1-t)>0&\quad\Longrightarrow\quad5\mid v_\ell(1-t),\\
v_\ell(t)<0&\quad\Longrightarrow\quad7\mid -v_\ell(t).
\end{align*}
When it exists, $\lambda$ is unique.
\end{theorem}
\begin{proof}
Write $m=v_\ell(t)$, $n=v_\ell(1-t)$ and $h=v_\ell(\lambda)$. Because $t+(1-t)=1$, the pair $(m,n)$ belongs to one of four sectors:
\[
(m>0,n=0),\quad(m=0,n>0),\quad(m=n<0),\quad(m=n=0).
\]
The output valuations are
\begin{align*}
a_\ell&=35h+12m+28n,\\
b_\ell&=21h+7m+17n,\\
c_\ell&=15h+5m+12n.
\end{align*}
They must be nonnegative, and their minimum must be zero. Equivalently, the minimum of $3a_\ell,5b_\ell,7c_\ell$ must be zero.

In the first sector these weighted valuations are $105h+36m,105h+35m,105h+35m$. Therefore $h=-m/3$, and the output is $(m/3,0,0)$. In the second sector they are $105h+84n,105h+85n,105h+84n$. This forces $h=-4n/5$ and output $(0,n/5,0)$. In the third sector write $m=n=-r$ with $r>0$. The weighted valuations are $105h-120r,105h-120r,105h-119r$. Thus $h=8r/7$ and the output is $(0,0,r/7)$. In the final sector $h=0$.

An integral $h$ exists in these four cases precisely under the stated divisibility conditions. Only finitely many primes occur in $t$ or $1-t$, so $\lambda=\prod_\ell\ell^{h_\ell}$ is a well defined positive rational number. The three output supports are disjoint, proving integrality and primitivity. Every $h_\ell$ was forced, so two positive choices of $\lambda$ would have ratio a positive rational with zero valuation at every prime. That ratio is one.
\end{proof}
\figplaceholder{03_valuation_packets}{The three nontrivial valuation sectors retain different divisibility conditions. A residue projector can preserve their labels without making the conditions hold.}
If $t=m/n$ is in lowest terms with $0<m<n$, the theorem says that $m$ must be a cube, $n-m$ a fifth power, and $n$ a seventh power. The parametrization has not escaped the original power equation. It has isolated its defect in three valuation sectors.

\begin{example}
At $t=1/2$, the denominator has valuation one at two, which fails the required divisibility by seven. No positive rational scaling produces a primitive integer triple. The choice $\lambda=4$ nevertheless gives
\[
(A,B,C)=(2^{30},2^{18},2^{13}),
\]
and the power equation holds. Its common divisor $2^{13}$ makes it a negative control against interpreting rational parametrization as a primitive solution.
\end{example}
This is where the Komorebi packet viewpoint has concrete force. Each sector carries an exact linear valuation rule, a required divisibility residue, and a reconstruction weight. Keeping all three records prevents information loss. The records alone do not prove that no rational $t$ satisfies every condition. A global closure still requires a relation linking the sectors, rather than additional names for the same local data.
